%======================================================================
\part*{Part I — Core-Guided and IHS MaxSAT as Benders and Dantzig--Wolfe Decomposition}
\addcontentsline{toc}{part}{Part I — Core-Guided and IHS MaxSAT as Benders and Dantzig--Wolfe Decomposition}
%======================================================================

\begin{abstract}
Maximum Satisfiability (MaxSAT) and Mixed-Integer Programming (MIP) are two
major frameworks for combinatorial optimization that have largely developed in
separate communities. In this part we draw explicit structural parallels
between two families of MaxSAT algorithms and two classical decomposition
methods from operations research. We show that core-guided MaxSAT algorithms
operate as instances of Logic-Based Benders Decomposition (LBBD), where
unsatisfiable cores play the role of Benders cuts derived via inference duality.
Dually, we argue that Implicit Hitting Set (IHS) MaxSAT solvers implement a
form of Dantzig--Wolfe decomposition when viewed through LP duality, with the
SAT oracle acting as a pricing subproblem that generates new constraints
(columns in the dual). These connections unify the algorithmic landscape across
the two communities and suggest avenues for cross-fertilization.
\end{abstract}

%----------------------------------------------------------------------
\section{Introduction}
\label{sec:intro}
%----------------------------------------------------------------------

Maximum Satisfiability (MaxSAT) is the optimization version of the Boolean
Satisfiability problem (SAT): given a set of hard clauses~$\mathcal{H}$ and a
set of weighted soft clauses~$(\mathcal{S}, w)$, find a truth assignment that
satisfies all hard clauses while maximizing the total weight of satisfied soft
clauses (equivalently, minimizing the cost of falsified soft clauses).

Modern exact MaxSAT solvers fall into two broad families:
\begin{enumerate}
\item \textbf{Core-guided algorithms} (originating with Fu and Malik~\citep{fu2006solving} and the unsatisfiable-core MSU family of~\citet{marquessilva2008algorithms}; later refinements include PMRES~\citep{narodytska2014maximum} and OLL~\citep{morgado2014core}), which iteratively call a SAT solver, extract unsatisfiable cores, and transform the formula to raise a lower bound.
\item \textbf{Implicit Hitting Set (IHS) algorithms} (e.g., MaxHS~\citep{davies2011solving,davies2013exploiting}), which alternate between a SAT oracle that extracts cores and an IP solver that computes a minimum-cost hitting set of the cores found so far.
\end{enumerate}

In operations research, Benders decomposition~\citep{benders1962partitioning}
and Dantzig--Wolfe decomposition~\citep{dantzigwolfe1960} are two classical
methods for structured optimization problems. Both decompose a hard monolithic
problem into a coordinating master problem and simpler subproblems, iterating
by passing cuts or columns between them. They are linked through LP
duality: Benders adds \emph{rows} (cuts) to a master problem defined over
complicating variables, while Dantzig--Wolfe adds \emph{columns} to a master
problem defined over complicating constraints.

The purpose of this part is to make explicit the structural correspondences:
\begin{itemize}
\item Core-guided MaxSAT $\longleftrightarrow$ Logic-Based Benders Decomposition,
\item IHS MaxSAT $\longleftrightarrow$ Dantzig--Wolfe Decomposition (in the dual).
\end{itemize}

%----------------------------------------------------------------------
\section{Preliminaries}
\label{sec:prelim}
%----------------------------------------------------------------------

\subsection{MaxSAT, Cores, and Correction Sets}

We fix notation for a weighted partial MaxSAT instance
$\langle \mathcal{H}, \mathcal{S}, w \rangle$ with hard clauses~$\mathcal{H}$,
soft clauses $\mathcal{S} = \{s_1,\ldots,s_m\}$, and weights
$w : \mathcal{S} \to \mathbb{Q}_{>0}$. A \emph{core} is a subset
$\kappa \subseteq \mathcal{S}$ with $\mathcal{H} \cup \kappa$ unsatisfiable; a
\emph{correction set} is a subset $\mathcal{C} \subseteq \mathcal{S}$ with
$\mathcal{H} \cup (\mathcal{S} \setminus \mathcal{C})$ satisfiable. By the
hitting-set duality of~\citet{liffiton2008algorithms}, every minimal core hits
every minimal correction set and vice versa---a fact we will use
repeatedly.

\subsection{Core-Guided and IHS Algorithms}

The two algorithm families differ in \emph{where} they accumulate information about the instance.

\paragraph{Core-guided algorithms.}
Core-guided solvers maintain the formula itself as an implicit representation of the lower bound. At each iteration a SAT oracle is called; if it returns unsatisfiable, a core~$\kappa$ is extracted and the formula is syntactically \emph{transformed} so that at least one soft clause in~$\kappa$ must be relaxed, raising the lower bound by the minimum weight in the core. Algorithm~\ref{alg:cg} gives a generic presentation that captures the Fu--Malik algorithm~\citep{fu2006solving}, its MSU-family extensions~\citep{marquessilva2008algorithms}, PMRES~\citep{narodytska2014maximum}, and OLL~\citep{morgado2014core} as instantiations of the \textsc{Relax} step.

\begin{algorithm}[t]
\caption{Generic Core-Guided MaxSAT}
\label{alg:cg}
\begin{algorithmic}[1]
\Require weighted partial MaxSAT instance $\langle \mathcal{H}, \mathcal{S}, w \rangle$
\Ensure optimal cost $\mathit{OPT}$ and a satisfying assignment
\State $\mathit{lb} \leftarrow 0$;\quad $\mathcal{S}' \leftarrow \mathcal{S}$;\quad $w' \leftarrow w$
\Loop
    \State $\mathit{result} \leftarrow \textsc{SAT}\!\left(\mathcal{H} \cup \mathcal{S}'\right)$
    \If{$\mathit{result} = \mathbf{SAT}$}
        \State \Return $(\sigma,\; \mathit{lb})$ \Comment{$\sigma$ is the satisfying assignment}
    \EndIf
    \State $\kappa \leftarrow \textsc{UnsatCore}\!\left(\mathcal{H},\mathcal{S}'\right)$ \Comment{$\kappa \subseteq \mathcal{S}'$, $\mathcal{H} \cup \kappa$ unsat}
    \State $\delta \leftarrow \min\{w'(s) : s \in \kappa\}$
    \State $\mathit{lb} \leftarrow \mathit{lb} + \delta$
    \State $(\mathcal{S}', w') \leftarrow \textsc{Relax}(\mathcal{S}', w', \kappa, \delta)$
    \Comment{add relaxation vars; reduce weights by $\delta$}
\EndLoop
\end{algorithmic}
\end{algorithm}

The \textsc{Relax} step introduces a fresh relaxation variable~$r_s$ for each $s \in \kappa$, weakening each clause from~$s$ to $s \vee r_s$, and adds a \emph{cardinality constraint} $\sum_{s \in \kappa} r_s \geq 1$ (or a soft version thereof weighted~$\delta$). Weights exceeding $\delta$ give rise to residual copies of the clause with weight $w'(s) - \delta$. Different choices of cardinality encoding yield the different named algorithms; the common invariant is that $\mathit{lb}$ is a valid lower bound on the optimal cost after every iteration.

\paragraph{Implicit Hitting Set algorithms.}
IHS solvers~\citep{davies2011solving,davies2013exploiting} separate the two responsibilities cleanly: an IP solver maintains an explicit \emph{master problem} that optimizes over the cores found so far, while the SAT oracle is used only to validate or refute the current candidate solution. Algorithm~\ref{alg:ihs} describes the MaxHS solver.

\begin{algorithm}[t]
\caption{Implicit Hitting Set MaxSAT (MaxHS)}
\label{alg:ihs}
\begin{algorithmic}[1]
\Require weighted partial MaxSAT instance $\langle \mathcal{H}, \mathcal{S}, w \rangle$
\Ensure optimal cost $\mathit{OPT}$ and a satisfying assignment
\State $\mathcal{K} \leftarrow \emptyset$ \Comment{set of discovered cores}
\Loop
    \State $H^* \leftarrow \textsc{SolveIP}\!\left(\min \textstyle\sum_{i} w(s_i)\,y_i \;:\; \sum_{s_i \in \kappa} y_i \geq 1\;\forall \kappa \in \mathcal{K},\; y \in \{0,1\}^{|\mathcal{S}|}\right)$
    \State $\mathit{result} \leftarrow \textsc{SAT}\!\left(\mathcal{H} \cup \left(\mathcal{S} \setminus H^*\right)\right)$
    \If{$\mathit{result} = \mathbf{SAT}$}
        \State \Return $\!\left(\sigma,\; \textstyle\sum_{s \in H^*} w(s)\right)$
        \Comment{$H^*$ is a minimum-cost correction set}
    \EndIf
    \State $\kappa \leftarrow \textsc{UnsatCore}\!\left(\mathcal{H},\, \mathcal{S} \setminus H^*\right)$
    \Comment{$\kappa$ is not hit by $H^*$}
    \State $\mathcal{K} \leftarrow \mathcal{K} \cup \{\kappa\}$
\EndLoop
\end{algorithmic}
\end{algorithm}

At termination of Algorithm~\ref{alg:ihs}, the optimal cost equals $\sum_{s \in H^*} w(s)$ because (i)~$H^*$ is a minimum hitting set of all known cores, giving a lower bound, and (ii)~the SAT oracle confirms that $\mathcal{S} \setminus H^*$ is satisfiable together with $\mathcal{H}$, so $H^*$ is an achievable correction set and hence an upper bound. The lower and upper bounds coincide, certifying optimality.

\subsection{The LP Relaxation of MaxSAT}

Introduce a binary variable $y_i \in \{0,1\}$ for each soft clause $s_i$,
where $y_i = 1$ means $s_i$ is falsified. The MaxSAT problem can be written as
the integer program:
\begin{align}
\min \quad & \sum_{i=1}^{m} w(s_i)\, y_i \label{eq:maxsat-obj}\\
\text{s.t.} \quad & \sum_{s_i \in \kappa} y_i \geq 1 \quad \forall \text{ cores } \kappa, \label{eq:maxsat-core}\\
& y_i \in \{0,1\}. \label{eq:maxsat-int}
\end{align}
Constraint~\eqref{eq:maxsat-core} states that every core must be ``hit'' by
at least one falsified soft clause. This is the \emph{hitting set}
formulation, and there are exponentially many core constraints. The LP
relaxation replaces~\eqref{eq:maxsat-int} with $0 \leq y_i \leq 1$.

\subsection{Benders Decomposition}

Recall that classical Benders~\citep{benders1962partitioning} fixes
complicating variables~$\bar{x}$ in the master, solves an LP subproblem
over~$y$, and derives optimality or feasibility cuts from the subproblem's LP
dual. \emph{Logic-Based Benders Decomposition} (LBBD)~\citep{hooker2003logic}
generalizes this: the subproblem can be any optimization or feasibility
problem, and cuts are derived from any sound \emph{inference dual}---a proof
of infeasibility or suboptimality---rather than from LP duality alone.
Algorithm~\ref{alg:benders} gives the LBBD template.

\begin{algorithm}[t]
\caption{Logic-Based Benders Decomposition (LBBD)}
\label{alg:benders}
\begin{algorithmic}[1]
\Require problem $\min\{c^{\top}x + v(x) : x \in X\}$ where $v(x)$ is the value of a subproblem
\Ensure optimal solution $(x^*, y^*)$ and cost $\mathit{OPT}$
\State $\mathit{UB} \leftarrow +\infty$;\quad $\mathcal{C} \leftarrow \emptyset$ \Comment{$\mathcal{C}$: accumulated Benders cuts}
\Loop
    \State $\bar{x} \leftarrow \arg\min\{c^{\top}x + z : x \in X,\; (x,z) \text{ satisfies all cuts in } \mathcal{C}\}$
    \Comment{solve master}
    \State $(z_\mathrm{MP}, \bar{x}) \leftarrow$ master objective and solution
    \If{$z_\mathrm{MP} \geq \mathit{UB}$} \textbf{break} \EndIf \Comment{lower bound meets upper bound}
    \State $\mathit{result} \leftarrow \textsc{SolveSubproblem}(\bar{x})$
    \If{$\mathit{result} = \mathbf{Feasible}$ \textbf{and} $c^{\top}\bar{x} + v(\bar{x}) < \mathit{UB}$}
        \State $\mathit{UB} \leftarrow c^{\top}\bar{x} + v(\bar{x})$;\quad $(x^*, y^*) \leftarrow (\bar{x},\, \bar{y})$
        \Comment{new incumbent}
    \EndIf
    \State $\phi \leftarrow \textsc{InferenceDual}(\bar{x}, \mathit{result})$
    \Comment{derive cut from subproblem proof}
    \State $\mathcal{C} \leftarrow \mathcal{C} \cup \{\phi\}$
\EndLoop
\State \Return $(x^*, y^*, \mathit{UB})$
\end{algorithmic}
\end{algorithm}

The key degree of freedom in LBBD is \textsc{InferenceDual}: when the subproblem
is an LP, it returns the classical Benders cut from dual multipliers; when the
subproblem is a SAT call, it returns a combinatorial cut (an unsatisfiable core).

\subsection{Dantzig--Wolfe Decomposition}

Dantzig--Wolfe~\citep{dantzigwolfe1960} dualizes the setting: it reformulates
an LP with complicating \emph{constraints} by representing the tractable
polyhedron $\mathcal{P} = \{x : Dx \leq d,\; x \geq 0\}$ via its extreme
points. A \emph{restricted master} works with a subset of these points
(columns), and a \emph{pricing subproblem} searches for one of negative reduced
cost. Algorithm~\ref{alg:dw} gives the column generation procedure.

\begin{algorithm}[t]
\caption{Dantzig--Wolfe Decomposition (Column Generation)}
\label{alg:dw}
\begin{algorithmic}[1]
\Require LP $\min\{c^{\top}x : Ax = b,\; x \in \mathcal{P}\}$ with $\mathcal{P}$ block-structured
\Ensure optimal value $z^*$ and solution $x^*$
\State Initialize restricted master with a feasible subset of columns $\Lambda \subseteq \{p^1,\ldots,p^K\}$
\Loop
    \State Solve restricted master LP; let $\pi$ be the dual variables for $Ax = b$
    \State $(\bar{c}, p^{\mathrm{new}}) \leftarrow \textsc{PricingSubproblem}(\pi)$
    \Comment{$\bar{c} = \min_{p \in \mathcal{P}}\,(c - A^{\top}\pi)^{\top}p$}
    \If{$\bar{c} \geq 0$}
        \State \Return $(z^*, x^*)$ \Comment{no improving column; current solution is optimal}
    \EndIf
    \State $\Lambda \leftarrow \Lambda \cup \{p^{\mathrm{new}}\}$
    \Comment{add column with negative reduced cost to master}
\EndLoop
\end{algorithmic}
\end{algorithm}

The restricted master LP provides a valid \emph{dual bound} (lower bound for a
minimization problem) at every iteration; this bound tightens monotonically as
columns are added and converges to the LP optimum when no negative-reduced-cost
column exists.

Through LP duality, Benders and Dantzig--Wolfe are intimately related: adding a
row (cut) to the Benders master corresponds to adding a column to the dual of
that master, which is precisely the Dantzig--Wolfe master. The two algorithms
are therefore dual views of the same iterative refinement process.

%----------------------------------------------------------------------
\section{Core-Guided MaxSAT as Benders Decomposition}
\label{sec:benders}
%----------------------------------------------------------------------

\subsection{The Decomposition Structure}

A core-guided MaxSAT algorithm operates in the following loop:
\begin{enumerate}
\item \textbf{Subproblem (SAT oracle):} Given the current (transformed) formula, call a SAT solver. If satisfiable, the model is optimal; stop. If unsatisfiable, obtain a core~$\kappa$.
\item \textbf{Master update:} Use the core~$\kappa$ to transform the formula (add relaxation variables, cardinality constraints) so that the lower bound increases.
\item Repeat.
\end{enumerate}

We now map this onto the LBBD framework of \citet{hooker2003logic}.

Throughout, $y_i \in \{0,1\}$ denotes the indicator that soft clause $s_i$ is
falsified, and $H^* = \{s_i : y_i = 1\}$ denotes the corresponding set of
falsified soft clauses. We use $H^*$ uniformly for the master's current
guess across both algorithm families: in IHS it is explicit and computed by an
IP solve; in core-guided algorithms it is implicit, induced by the relaxation
variables introduced so far.

\begin{center}
\begin{tabular}{lll}
\toprule
\textbf{LBBD component} & \textbf{Core-guided MaxSAT analog} \\
\midrule
Complicating variables $x$ & Falsification indicators $y_i$ (set $H^*$) \\
Master problem & Accumulated lower bound + cardinality constraints \\
Subproblem $\mathrm{SP}(\bar{x})$ & SAT call on $\mathcal{H} \cup \mathcal{S}'$ \\
Infeasibility certificate & Unsatisfiable core $\kappa$ \\
Benders cut & Hitting set constraint $\sum_{s_i \in \kappa} y_i \geq 1$ \\
Inference dual & SAT refutation proof (e.g., resolution) \\
\bottomrule
\end{tabular}
\end{center}

\subsection{The Inference Dual Perspective}

In classical Benders, the cut $\pi^T b \leq z$ is derived from the LP dual of
the subproblem. In LBBD, the cut comes from any valid proof. In core-guided
MaxSAT, the SAT solver provides a refutation proof (via conflict-driven clause
learning) that establishes the unsatisfiability of $\mathcal{H} \cup \kappa$.
This proof serves as the solution to the inference dual: it certifies that
\emph{at least one} soft clause in $\kappa$ must be falsified, producing the
cut
\[
\sum_{s_i \in \kappa} y_i \geq 1.
\]
This is a \emph{combinatorial Benders cut}~\citep{hooker2003logic,codato2006combinatorial}: a linear inequality derived from a logical deduction rather than from LP duality.

\subsection{Formula Transformation as Implicit Cut Accumulation}

Core-guided solvers such as PMRES and OLL do not maintain an explicit master
IP. Instead, they transform the formula by adding relaxation variables and
cardinality constraints. Recent work by \citet{katsirelos2023analysis} has
shown that these transformations implicitly encode an LP. Specifically:

\begin{proposition}[\citeauthor{katsirelos2023analysis}, 2023]
The transformations performed by PMRES and OLL implicitly discover cores of the
original formula. In unweighted instances, the relaxed formula at each
iteration encodes a minimum hitting set of these cores. Moreover, the process
can be captured by a compact integer linear program whose LP relaxation yields
the same bound.
\end{proposition}

Thus, what appears to be a purely syntactic formula transformation is, in
fact, an implicit construction of the Benders master problem, with each core
extraction step corresponding to one iteration of the LBBD loop.

\begin{figure}[t]
\centering
\begin{tikzpicture}[
    node distance=2.5cm and 4cm,
    block/.style={draw, rounded corners, minimum width=3cm, minimum height=1.2cm, align=center, thick},
    arrow/.style={-{Stealth[length=3mm]}, thick},
]
    \node[block] (master) {Master\\(lower bound +\\accumulated cuts)};
    \node[block, right=of master] (sub) {Subproblem\\(SAT oracle)};

    \draw[arrow] (master.20) -- node[above,font=\small] {transformed formula $\mathcal{H} \cup \mathcal{S}'$} (sub.160);
    \draw[arrow] (sub.200) -- node[below,font=\small] {core $\kappa$ (Benders cut)} (master.340);

    \node[below=0.5cm of master,font=\small\itshape,align=center] {Encode falsified set\\implicitly via relaxation};
    \node[below=0.5cm of sub,font=\small\itshape,align=center] {Check feasibility\\of remaining formula};
\end{tikzpicture}
\caption{Core-guided MaxSAT as Logic-Based Benders Decomposition. The SAT oracle acts as the subproblem; each unsatisfiable core yields a combinatorial Benders cut.}
\label{fig:benders}
\end{figure}

%----------------------------------------------------------------------
\section{IHS MaxSAT as Dantzig--Wolfe Decomposition}
\label{sec:dw}
%----------------------------------------------------------------------

\subsection{The Implicit Hitting Set Framework}

The IHS algorithm~\citep{davies2011solving,moreno2013implicit} operates as follows:
\begin{enumerate}
\item \textbf{Master problem (IP solver):} Compute a minimum-cost hitting set $H^*$ of all cores found so far. That is, solve
\[
\min \sum_{i=1}^{m} w(s_i)\, y_i \quad \text{s.t.} \quad \sum_{s_i \in \kappa_j} y_i \geq 1 \;\; \forall j = 1,\ldots,k, \quad y_i \in \{0,1\}.
\]
\item \textbf{Oracle (SAT solver):} Check whether removing the clauses in $H^*$ from $\mathcal{S}$ makes the formula satisfiable. If yes, $H^*$ is optimal; stop. If not, the SAT solver returns a new core $\kappa_{k+1}$ that is not hit by $H^*$.
\item Add $\kappa_{k+1}$ as a new constraint and repeat.
\end{enumerate}

\subsection{The Dantzig--Wolfe Connection via LP Duality}

To see the Dantzig--Wolfe structure, consider the LP dual of the hitting set
formulation~\eqref{eq:maxsat-obj}--\eqref{eq:maxsat-core}. Let $\mathcal{K} =
\{\kappa_1, \ldots, \kappa_N\}$ be the set of all cores, and let $\lambda_j
\geq 0$ be the dual variable for core $\kappa_j$. The dual LP is:
\begin{align}
\max \quad & \sum_{j=1}^{N} \lambda_j \label{eq:dual-obj}\\
\text{s.t.} \quad & \sum_{j \,:\, s_i \in \kappa_j} \lambda_j \leq w(s_i) \quad \forall i = 1,\ldots,m, \label{eq:dual-pack}\\
& \lambda_j \geq 0. \label{eq:dual-nn}
\end{align}
This is a \emph{packing} problem over cores: find a maximum-weight fractional
packing of cores subject to the capacity of each soft clause.

Since the number of cores~$N$ is exponential, this dual cannot be solved
directly. However, it has the structure required for \emph{column generation}
(the algorithmic engine of Dantzig--Wolfe decomposition):
\begin{itemize}
\item The \textbf{restricted master} is the dual LP~\eqref{eq:dual-obj}--\eqref{eq:dual-nn} restricted to the cores $\kappa_1, \ldots, \kappa_k$ discovered so far. Each core $\kappa_j$ corresponds to a \emph{column} (variable $\lambda_j$) in this dual.
\item The \textbf{pricing subproblem} asks: does there exist a core $\kappa$ not in the current restricted master that has positive reduced cost? In the dual~\eqref{eq:dual-obj}--\eqref{eq:dual-nn}, this amounts to finding a core $\kappa$ such that $\mathcal{H} \cup \kappa$ is unsatisfiable and $\kappa$ is not hit by the current primal solution. The SAT oracle answers exactly this question.
\end{itemize}

\begin{center}
\begin{tabular}{lll}
\toprule
\textbf{Dantzig--Wolfe component} & \textbf{IHS MaxSAT analog} \\
\midrule
Restricted master (dual) & Packing LP over known cores \\
Column & Core $\kappa_j$ (dual variable $\lambda_j$) \\
Pricing subproblem & SAT oracle: find unhit core \\
Negative reduced cost & Core not hit by current hitting set \\
\midrule
Restricted master (primal) & Hitting set IP over known cores \\
New constraint (row) & New core $\kappa_{k+1}$ \\
\bottomrule
\end{tabular}
\end{center}

\begin{proposition}
The IHS algorithm is a primal-side implementation of column generation on
the dual of the MaxSAT hitting set LP. Each new core discovered by the SAT
oracle acts simultaneously as:
\begin{enumerate}
\item a new \emph{row} (constraint) in the primal hitting set master, and
\item a new \emph{column} (variable) in the dual packing master.
\end{enumerate}
By LP duality, adding a row to the primal is equivalent to adding a column to
the dual. Hence the IHS algorithm is a Dantzig--Wolfe procedure on the dual
formulation.
\end{proposition}

\begin{figure}[t]
\centering
\begin{tikzpicture}[
    node distance=2.5cm and 4cm,
    block/.style={draw, rounded corners, minimum width=3cm, minimum height=1.2cm, align=center, thick},
    arrow/.style={-{Stealth[length=3mm]}, thick},
]
    \node[block] (master) {Master\\(hitting set IP /\\dual packing LP)};
    \node[block, right=of master] (sub) {Pricing subproblem\\(SAT oracle)};

    \draw[arrow] (master.20) -- node[above,font=\small] {hitting set $H^*$} (sub.160);
    \draw[arrow] (sub.200) -- node[below,font=\small,align=center] {new core $\kappa$\\(row / dual column)} (master.340);

    \node[below=0.5cm of master,font=\small\itshape,align=center] {Solve restricted\\hitting set};
    \node[below=0.5cm of sub,font=\small\itshape,align=center] {Find core not\\hit by $H^*$};
\end{tikzpicture}
\caption{IHS MaxSAT as Dantzig--Wolfe Decomposition (dual view). The SAT oracle acts as the pricing subproblem; each new core is a column in the dual packing LP.}
\label{fig:dw}
\end{figure}

\subsection{LP Bounds and Convergence}

In Dantzig--Wolfe decomposition, the restricted master LP provides a bound
that improves monotonically as columns are added, converging to the LP
optimum. Similarly, in IHS, adding more cores tightens the LP relaxation of
the hitting set master. \citet{davies2013exploiting} showed that solving the
hitting set master as an IP (rather than an LP) accelerates convergence
in practice by providing both primal feasible solutions and dual bounds
simultaneously.

The work of \citet{katsirelos2023analysis} further showed that core-guided
algorithms such as OLL implicitly compute feasible dual solutions of an LP
that encodes the same hitting set structure. This means that both families
of algorithms---core-guided and IHS---are optimizing over the \emph{same}
LP relaxation, but from opposite sides of duality.

\subsection{What the Analogy Preserves, and What It Doesn't}
\label{sec:preserves}

The Dantzig--Wolfe view is illuminating but is not a textbook instantiation
of column generation. To use it responsibly, it helps to be explicit about
which structural features carry over and which do not.

\paragraph{What the analogy preserves.} The most robust part of the analogy
is the \emph{row--column duality between cuts and cores}. Adding a hitting-set
row $\sum_{s_i \in \kappa} y_i \geq 1$ to the primal IP
formulation~\eqref{eq:maxsat-obj}--\eqref{eq:maxsat-int} is exactly equivalent,
by LP duality, to introducing a new variable~$\lambda_\kappa$ in the dual
packing LP~\eqref{eq:dual-obj}--\eqref{eq:dual-nn}. This is the same
row-to-column transformation that mediates the Benders/Dantzig--Wolfe duality
in classical OR. Likewise, the \emph{reduced-cost interpretation of an unhit
core} is exact: given the current dual prices~$\bar{y}$ (a fractional hitting
set from the LP relaxation of the master), a core $\kappa$ has positive
reduced cost in the dual packing problem iff
$\sum_{s_i \in \kappa} \bar{y}_i < 1$, which is precisely the condition that
$\kappa$ is not (fractionally) hit. The SAT oracle is therefore a sound
implementation of pricing: if it finds a new core, the master's LP bound is
not yet optimal; if it cannot, the master's solution is LP-optimal.

\paragraph{What the analogy does not preserve.} Three classical features of
Dantzig--Wolfe break down here, and recognizing this prevents over-claiming.

\begin{enumerate}
\item \textbf{Pricing is NP-complete, not polynomial.} In textbook
Dantzig--Wolfe the subproblem is a tractable LP with block structure (network
flow, knapsack, etc.) and pricing returns a column in polynomial time. Here
pricing is a SAT call---NP-complete in general. The high per-iteration cost
is real, and many of the convergence-rate analyses for column generation (e.g.,
worst-case iteration bounds in terms of LP geometry) do not transfer directly.
What does transfer is that pricing returns a column with negative reduced cost
\emph{when one exists}; modern SAT solvers are sound and complete oracles,
which is what the analogy actually requires.
\item \textbf{The master is solved as an IP, not an LP.} Practical IHS solves
the master as an integer program to obtain feasible hitting sets directly
(\citealp{davies2013exploiting}), which is non-standard for column generation
and closer to branch-and-price without the branching. In particular, the
\emph{LP relaxation of the master is not what is solved at each iteration},
so the dual bound used by classical pricing arguments is not directly
available. The analogy holds at the level of the underlying LP, but the
implementation departs from it.
\item \textbf{Columns are generated, not enumerated convex combinations.}
Standard Dantzig--Wolfe represents the subproblem polytope $\mathcal{P}$ via
its extreme points, with the master expressing $x$ as a convex combination
$x = \sum_k \lambda_k p^k$. In MaxSAT there is no such polytope to
parameterize; cores are constraints rather than points, and the columns
$\lambda_\kappa$ live only in the \emph{dual}. Equivalently: the
MaxSAT version of Dantzig--Wolfe is column generation on the
\emph{constraint-generation dual} of an exponentially-constrained covering
problem, not on a Minkowski decomposition of a primal polytope. This is a
known but more specialized setting in OR (sometimes called ``constraint
generation viewed dually as column generation''), and it is the version
that fits MaxSAT.
\end{enumerate}

\paragraph{Net assessment.} The analogy is sound at the level of LP duality
between primal cuts and dual columns, and at the level of the pricing oracle's
specification. It is loose at the level of complexity, convergence rates, and
master tractability. The useful consequence is that techniques whose
correctness depends only on the row--column duality---in-out separation,
dual stabilization, Lagrangian smoothing of master prices---should transfer.
Techniques that rely on polynomial pricing or LP master geometry should not
be expected to transfer without adaptation.

%----------------------------------------------------------------------
\section{A Worked Example}
\label{sec:example}
%----------------------------------------------------------------------

To make the row--column duality concrete, we trace both algorithms on the
same small instance and compare what each one accumulates.

\paragraph{Instance.} No hard clauses; four unit soft clauses, all of
weight~1:
\[
  s_1 = a,\quad s_2 = \neg a,\quad s_3 = b,\quad s_4 = \neg b.
\]
The optimum is $\mathit{OPT} = 2$: any solution must falsify exactly one of
$\{s_1, s_2\}$ and one of $\{s_3, s_4\}$. The only minimal cores are
$\kappa_1 = \{s_1, s_2\}$ and $\kappa_2 = \{s_3, s_4\}$.

\paragraph{Hitting-set LP.} The MaxSAT IP and its LP relaxation are
\[
  \min\; y_1+y_2+y_3+y_4
  \;\;\text{s.t.}\;\; y_1+y_2 \geq 1,\;\; y_3+y_4 \geq 1,\;\; y \in [0,1]^4,
\]
with LP value~$2$ (integrality gap zero, since this matrix is balanced). The
dual packing LP is
\[
  \max\; \lambda_1 + \lambda_2
  \;\;\text{s.t.}\;\; \lambda_1 \leq w(s_1),\, \lambda_1 \leq w(s_2),\,
  \lambda_2 \leq w(s_3),\, \lambda_2 \leq w(s_4),\, \lambda \geq 0,
\]
with optimum $\lambda_1 = \lambda_2 = 1$, value~$2$.

\paragraph{Trace of core-guided (Algorithm~\ref{alg:cg}).}
\begin{enumerate}
\item $\textsc{SAT}(\mathcal{H} \cup \mathcal{S}')$ on the four units:
\textbf{UNSAT}. Return core $\kappa_1 = \{s_1, s_2\}$. Set
$\delta = 1$, $\mathit{lb} = 1$. \textsc{Relax} weakens
$s_1 \mapsto a \vee r_1$ and $s_2 \mapsto \neg a \vee r_2$ and adds the soft
cardinality constraint $r_1 + r_2 \geq 1$ (weight~1).
\item $\textsc{SAT}$ on the relaxed formula: \textbf{UNSAT}. The new core
is $\kappa_2 = \{s_3, s_4\}$. Set $\delta = 1$, $\mathit{lb} = 2$. Relax
$s_3, s_4$ analogously, adding $r_3 + r_4 \geq 1$.
\item $\textsc{SAT}$: \textbf{SAT}. Set, e.g., $a = \mathit{true}$,
$b = \mathit{true}$, $r_2 = r_4 = 1$, $r_1 = r_3 = 0$. Return cost
$\mathit{lb} = 2$.
\end{enumerate}

\paragraph{Trace of IHS (Algorithm~\ref{alg:ihs}).}
\begin{enumerate}
\item $\mathcal{K} = \emptyset$. The IP master is unconstrained, so $H^* =
\emptyset$. $\textsc{SAT}$ on the full formula: \textbf{UNSAT}, core
$\kappa_1 = \{s_1, s_2\}$.
\item $\mathcal{K} = \{\kappa_1\}$. The IP master $\min\sum y_i$ s.t.\
$y_1 + y_2 \geq 1$ gives, e.g., $H^* = \{s_1\}$ (cost~1). $\textsc{SAT}$ on
$\{s_2, s_3, s_4\}$: \textbf{UNSAT}, core $\kappa_2 = \{s_3, s_4\}$.
\item $\mathcal{K} = \{\kappa_1, \kappa_2\}$. The IP master gives, e.g.,
$H^* = \{s_1, s_3\}$ (cost~2). $\textsc{SAT}$ on $\{s_2, s_4\}$ = $\neg a
\wedge \neg b$: \textbf{SAT}. Return cost~$2$.
\end{enumerate}

\paragraph{Side-by-side comparison.} Both algorithms discover exactly the
same two cores $\{\kappa_1, \kappa_2\}$, in the same order, and terminate with
the same optimum. What differs is the data structure that accumulates them
(Table~\ref{tab:trace}):

\begin{table}[h]
\centering
\caption{The two algorithms accumulate the same information in dual forms.}
\label{tab:trace}
\begin{tabular}{lll}
\toprule
\textbf{Iteration} & \textbf{Core-guided accumulates} & \textbf{IHS accumulates} \\
\midrule
1 ($\kappa_1$ found) & relaxation vars $r_1, r_2$;     & primal row $y_1 + y_2 \geq 1$;     \\
                     & cardinality $r_1 + r_2 \geq 1$  & dual column $\lambda_1$            \\
2 ($\kappa_2$ found) & relaxation vars $r_3, r_4$;     & primal row $y_3 + y_4 \geq 1$;     \\
                     & cardinality $r_3 + r_4 \geq 1$  & dual column $\lambda_2$            \\
\midrule
After 2 iter.\ & $\mathit{lb} = 2$ encoded by         & $H^* \in \{ \{s_1,s_3\}, \{s_1,s_4\}, \{s_2,s_3\}, \{s_2,s_4\} \}$, \\
               & relaxation structure                  & all of cost~2 \\
\bottomrule
\end{tabular}
\end{table}

The two cardinality constraints accumulated on the core-guided side are
exactly the two hitting-set rows on the IHS side---syntactically rewritten on
the relaxation variables rather than the original $y_i$, but encoding the
same information. By LP duality, those two rows correspond to two dual
columns $\lambda_1, \lambda_2$ in the packing LP. After two iterations both
sides have identified the LP optimum~$\lambda_1 = \lambda_2 = 1 = \mathit{OPT}$,
one by raising $\mathit{lb}$ through formula transformations (Benders-style),
the other by computing it explicitly via an IP master (Dantzig--Wolfe-style on
the dual).

%----------------------------------------------------------------------
\section{The Unified Picture}
\label{sec:unified}
%----------------------------------------------------------------------

The preceding analysis reveals a clean duality diagram connecting the four
methods:

\begin{figure}[h]
\centering
\begin{tikzpicture}[
    node distance=3cm and 5cm,
    block/.style={draw, rounded corners, minimum width=3.5cm, minimum height=1cm, align=center, thick},
    arrow/.style={{Stealth[length=3mm]}-{Stealth[length=3mm]}, thick, dashed},
]
    \node[block] (benders) {Benders\\Decomposition};
    \node[block, right=of benders] (dw) {Dantzig--Wolfe\\Decomposition};
    \node[block, below=of benders] (cg) {Core-Guided\\MaxSAT};
    \node[block, below=of dw] (ihs) {IHS\\MaxSAT};

    \draw[arrow] (benders) -- node[above,font=\small] {LP duality} (dw);
    \draw[arrow] (cg) -- node[below,font=\small] {LP duality} (ihs);
    \draw[arrow] (benders) -- node[left,font=\small,align=right] {inference\\dual} (cg);
    \draw[arrow] (dw) -- node[right,font=\small,align=left] {column gen.\\= core gen.} (ihs);
\end{tikzpicture}
\caption{The duality square connecting decomposition methods and MaxSAT algorithms.}
\label{fig:square}
\end{figure}

\begin{itemize}
\item \textbf{Horizontal axis (LP duality):} Benders and Dantzig--Wolfe are
dual to each other; core-guided and IHS are dual to each other in the same
sense. Adding a cut in one is adding a column in the other.
\item \textbf{Vertical axis (generalization):} Core-guided MaxSAT is an
instance of LBBD where the inference dual replaces the LP dual. IHS MaxSAT
is an instance of Dantzig--Wolfe where the pricing oracle is a SAT solver
rather than an LP.
\end{itemize}

This diagram also explains empirically observed complementarities. IHS solvers
(like Dantzig--Wolfe) tend to produce good primal solutions quickly because
the master IP yields feasible hitting sets at each iteration. Core-guided
solvers (like Benders) tend to produce strong dual bounds because each core
directly tightens the relaxation. This mirrors the well-known tradeoff in
OR between cut-based and column-based approaches.

%----------------------------------------------------------------------
\section{Transferable Techniques from OR to MaxSAT}
\label{sec:techniques}
%----------------------------------------------------------------------

Section~\ref{sec:preserves} identified exactly which structural features carry
over from Dantzig--Wolfe and Benders to the MaxSAT setting. The useful
consequence is that techniques whose correctness depends only on row--column
duality and the soundness of the pricing oracle should transfer essentially
unchanged. This section collects the most actionable candidates, in roughly
increasing order of implementation cost.

\subsection{Cross-Family Warm-Starting}

The cheapest hybrid technique is initialization. A short run of one algorithm
produces information that warm-starts the other, just as OR practitioners
routinely warm-start Benders or Dantzig--Wolfe masters with cuts or columns
from a heuristic phase.

\begin{itemize}
\item \emph{IHS seeded by core-guided.} Run Algorithm~\ref{alg:cg} for a fixed
iteration budget; harvest all cores it has discovered (whether explicitly
extracted or implicit in the relaxation structure). Initialize
Algorithm~\ref{alg:ihs} with $\mathcal{K}$ set to this collection. The first
IP master solve already returns a non-trivial $H^*$, skipping the cold-start
phase where IHS has no master structure to exploit.
\item \emph{Core-guided seeded by IHS or MCS enumeration.} Conversely, a quick
enumeration of minimal correction sets~\citep{liffiton2008algorithms} yields
cores by hitting-set duality. Pre-relaxing the formula along these cores gives
Algorithm~\ref{alg:cg} a head start on the lower bound.
\end{itemize}

This requires no algorithmic surgery---only sharing the data structure that
holds discovered cores---and so is the most accessible of the techniques
listed.

\subsection{Pareto-Optimal Core Selection}

A SAT oracle that returns UNSAT typically admits many valid cores; modern
conflict-driven solvers return one based on internal heuristics, not on what
is most useful to the master. The Benders literature provides a principled
answer: \citet{magnanti1981accelerating} select, among multiple valid cuts,
the one that is \emph{non-dominated} (Pareto-optimal) with respect to a
reference interior point of the master's feasible region.

The MaxSAT analog is concrete. Given an UNSAT call and a current fractional
hitting set~$\bar{y}$ (the LP relaxation of the master), prefer cores
$\kappa$ that maximize the violation $1 - \sum_{s_i \in \kappa} \bar{y}_i$.
Among cores with the same violation, prefer those that are minimal. This
couples MUS extraction with the master's current state, replacing the SAT
solver's internal heuristic with a Magnanti--Wong-style criterion. Because
soundness does not depend on which specific core is returned, this is a pure
orthogonal improvement and can be layered on top of existing IHS or
core-guided implementations.

\subsection{Dual Stabilization for IHS}

The IHS master, viewed as the dual of the packing
LP~\eqref{eq:dual-obj}--\eqref{eq:dual-nn}, suffers from the same instabilities
that plague column generation: degenerate dual solutions cause $H^*$ to
oscillate between alternative optima, slowing convergence in the late phase
(``tailing off''). Three OR techniques address this directly.

\paragraph{Wentges-style box stabilization.}
\citet{wentges1997weighted} adds a penalty term to the master that discourages
dual prices from drifting too far from a current \emph{stability center}.
Translated to IHS: maintain a reference fractional hitting set~$\bar{y}$ and
replace the master objective $\sum_i w(s_i)\, y_i$ by
$\sum_i w(s_i)\, y_i + \alpha \|y - \bar{y}\|_1$ for a decaying
penalty~$\alpha$. The stability center is updated when the SAT oracle
confirms progress. Implementation cost is low (one additional term in the IP
master), and correctness is preserved because $\alpha \to 0$ at termination.

\paragraph{In-out separation.}
\citet{benameur2007stabilization} interpolate between the current dual
solution and an interior point of the feasible dual region before pricing.
The MaxSAT analog: instead of calling the SAT oracle on $\mathcal{S}
\setminus H^*$, call it on $\mathcal{S} \setminus H^{\mathrm{io}}$, where
$H^{\mathrm{io}}$ is a convex combination of $H^*$ with a known feasible
(possibly suboptimal) hitting set. Cores extracted from interior points
tend to be deeper---they hit more constraints simultaneously---and they
damp master oscillation.

\paragraph{Bundle / proximal methods.}
For large-scale instances, bundle methods maintain a polyhedral approximation
of the dual function with proximal regularization. They could be adapted to
IHS by keeping a small set of past hitting sets and choosing the next probe
via a proximal step; this is more disruptive to implement but is the standard
heavy-machinery response to tailing-off in column generation.

\subsection{Price-Aware Core Extraction}

A more invasive idea draws on Lagrangian relaxation. The dual prices~$\bar{y}$
from the master's LP relaxation can be interpreted as a measure of how
``expensive'' each soft clause is to falsify. Cores composed of high-price
soft clauses contribute more to tightening the bound (a higher minimum
weight~$\delta$ in Algorithm~\ref{alg:cg}; a tighter packing constraint in
the dual~\eqref{eq:dual-pack}). One can therefore bias the SAT oracle's
search toward such cores---e.g., by reweighting VSIDS scores by current dual
prices, or by augmenting branching heuristics with the master's dual
information.

This is the most speculative of the techniques listed: it requires
SAT-solver modification, not just outer-loop changes. But it has a clean
theoretical foundation: it is the MaxSAT analog of running the pricing
subproblem with the current dual prices as objective, which is the defining
operation of column generation.

\subsection{Toward Hybrid Primal--Dual Solvers}

Together, the techniques above point toward a unified solver that maintains
both representations simultaneously: an explicit IP master (IHS-side) and a
transformed working formula (core-guided side). At each iteration, the
algorithm chooses whether the next core should be extracted by an IHS-style
probe (using $H^*$ as assumptions) or by a core-guided probe (on the current
relaxed formula), based on which is more likely to produce a useful cut. The
LP-based reformulation of OLL by~\citet{katsirelos2025core} provides a
natural starting point, since it makes the dual computations underlying
core-guided algorithms explicit and therefore composable with an IHS-style
IP master.

%----------------------------------------------------------------------
\section{Discussion}
\label{sec:discussion}
%----------------------------------------------------------------------

\paragraph{Implications for algorithm design.}
The decomposition viewpoint suggests a portfolio of techniques transferable
from OR to MaxSAT; these are developed in Section~\ref{sec:techniques}, from
the most readily deployable (cross-family warm-starting) to the most
speculative (price-aware core extraction inside the SAT solver).

\paragraph{Limitations.}
The analogy is structural, not exact. Section~\ref{sec:preserves} gives the
detailed breakdown of what is preserved (row--column duality between cuts and
cores; soundness of the SAT oracle as a pricing routine) and what is not
(polynomial pricing, LP master geometry, primal Minkowski decomposition).

%----------------------------------------------------------------------
\section{Future Work}
\label{sec:future}
%----------------------------------------------------------------------

The viewpoint developed here makes several concrete predictions and surfaces
open questions worth pursuing.

\paragraph{Empirical validation of transferred techniques.} The techniques
collected in Section~\ref{sec:techniques} are presented on theoretical
grounds. Each one yields a falsifiable empirical claim: cross-family
warm-starting should reduce IHS cold-start iterations on benchmark families
where core-guided is known to find good lower bounds quickly; Wentges-style
stabilization should reduce the tailing-off plateau in IHS solve time on
instances with degenerate hitting-set LPs; Pareto-optimal core selection
should reduce iteration counts on instances with many redundant cores.
Systematic experimental study of these predictions---ideally on the MSE
benchmarks---would be the most direct way to validate (or refute) the
operational value of the analogy. We are not aware of any such study to date.

\paragraph{The weighted case.} The cleanest LP encodings of core-guided
algorithms~\citep{katsirelos2023analysis} are for unweighted instances. The
weighted case is where most practical interest lies, and where the
bound-tightening story is more delicate: weight-stratified core extraction,
weight splitting, and the interaction with $\delta$ in
Algorithm~\ref{alg:cg} change the structure of the implicit master in ways
that the unweighted analysis does not directly cover. Establishing whether the
Dantzig--Wolfe analogy is exact in the weighted case, or merely heuristic, is
an open question.

\paragraph{Filling in the square.} Figure~\ref{fig:square} has four corners,
but OR offers richer machinery built on top of each: \emph{branch-and-price}
extends Dantzig--Wolfe with branching on master variables; \emph{Benders
branch-and-cut} interleaves cut generation with branching in a single
search tree. The MaxSAT analogs of these are largely missing from the
landscape. IHS with branching on hitting-set decisions is one natural
direction; core-guided algorithms equipped with an explicit IP master in the
spirit of~\citet{davies2013exploiting} are arguably a partial Benders
branch-and-cut. Naming the empty cells of the analogy is often where the
genuine research opportunities lie.

%----------------------------------------------------------------------
\section{Part I Conclusion}
\label{sec:conclusion}
%----------------------------------------------------------------------

We have presented explicit structural correspondences between two families of
MaxSAT algorithms and two classical decomposition methods from operations
research. Core-guided MaxSAT algorithms implement a form of Logic-Based
Benders Decomposition, where unsatisfiable cores serve as combinatorial
Benders cuts derived through inference duality. Implicit Hitting Set
algorithms implement a form of Dantzig--Wolfe decomposition (viewed through LP
duality), where the SAT oracle acts as a pricing subproblem generating new
columns in the dual packing formulation. These connections are mediated by LP
duality: the hitting set (primal) and core packing (dual) formulations of
MaxSAT stand in the same relationship as Benders and Dantzig--Wolfe masters.

Making these connections explicit opens a two-way transfer of techniques
between the OR and SAT communities, from cut strengthening and stabilization
to hybrid primal-dual algorithms that leverage the strengths of both
decomposition paradigms.


